\section{Where ATFM fits, and what would make it useful}
\label{sec:atfm}
\subsection{The specific information advantage being proposed}
ATFM observes session/tool events, maintains a live forecast board, and turns
forecasts into admission or placement-related directives. Its potential
advantage is \emph{lead time}: a tool-progress update may reveal a return before
a gateway sees the next LLM request. A router choosing among destinations for
an already-arrived prompt solves a different timing problem, although modern
routing and scaling systems already use predictive models.

\fig{atfm}{0.97}{ATFM's conceptual feedback loop. Forecasts become bounded
control decisions; observed outcomes must validate that decisions were delivered
and useful. Cache actuation is backend-dependent and is not depicted as a
universal direct GPU pin.}{atfm}

The candidate contribution has three parts: conditional forecasts from
session-specific observations; aggregate resource uncertainty rather than only
mean request rate; and decisions made early enough to affect future contention.
Each must earn its complexity. A high-quality forecast delivered after the
request arrives has no prefetch lead-time advantage. A prediction that changes
no action has no decision value. An action that increases another tenant's
latency more than it saves may be undesirable even when the forecast is correct.

The architecture suggests a useful separation. Keep the request path bounded
and capable of falling back. Run heavier forecasting on snapshots or an
independently budgeted path. Deliver versioned, expiring decisions to actuators
that own the affected resource. This is a design direction; it must not be
mistaken for a claim that every isolation and acknowledgement mechanism is
already implemented.

\subsection{Admission constraints: population risk and sample risk}
For resource $r$ in slot $s$, let $Y_{rs}$ be uncertain protected demand,
$a_{is}$ a binary assignment of background job $i$, and $d_{ir}$ its estimated
demand. A conceptual marginal chance constraint is
\begin{equation}
 \Prb\!\left(Y_{rs}+\sum_i a_{is}d_{ir}\le C_{rs}\right)
 \ge1-\epsilon_{rs}.
 \label{eq:chance}
\end{equation}
Capacity $C_{rs}$ must use the same units and time interpretation as demand.
Existing resident KV and future arriving KV are not interchangeable. A system
that imposes separate resource constraints does not automatically have a joint
$1-\epsilon$ guarantee across resources and slots. A conservative union bound
allocates risks so $\sum_{r,s}\epsilon_{rs}\le\epsilon_{\mathrm{joint}}$;
a joint-scenario method is another option.

With $K$ forecast samples, a planner can evaluate
\begin{equation}
 \frac1K\sum_{k=1}^K\ind\left\{Y_{rs}^{(k)}+\sum_i a_{is}d_{ir}\le C_{rs}\right\}
 \ge1-\epsilon_{rs}.
\end{equation}
This is an empirical constraint under the forecast sample distribution. It is
not automatically a frequentist guarantee for true future traffic. Finite
sample uncertainty, model error, drift, and correlated residuals all matter.
As a simple illustration, even for independent correctly generated draws, a
Bernoulli event with probability 0.1 has estimated-probability standard error
$\sqrt{0.1\cdot0.9/256}\approx0.0188$ at 256 samples. Increasing samples reduces
Monte Carlo noise but does not correct model bias.

ATFM's GDP planner assigns deferrable work in forecast order using sampled
resource thresholds, bounded holds, and tenant-delay limits. Delay caps can
release work even when the capacity test has no feasible slot. Consequently
this implementation should be described as a bounded-delay heuristic with
empirical feasibility checks, not a proof of overload prevention. The code and
its scaling notes state the mechanics; a deployable capacity model remains an
experimental responsibility.

\subsection{A worked use case and its failure cases}
Imagine a worker pool near a latency limit. Several background agents are
running tests; an interactive user is generating a response. Some agents report
that their tests are nearly complete and will soon return with long contexts.
If this information is reliable and arrives early, a controller could defer
new deferrable tool launches, preserve selected cache state, or prepare capacity.
The objective is to reduce the coming peak, not to speed up the tests or invent
more GPU capacity.

Now change the assumptions. If the tests emit misleading progress, the
controller can hold unrelated work unnecessarily. If the contexts are small,
prefetch/control overhead can exceed the saved prefill. If the engine has ample
headroom, delaying anything may only increase completion time. If link bandwidth
is saturated, prefetching the predicted burst can congest it earlier. If an
existing TTL policy already retains the right prefixes, the extra forecast may
be redundant. If tool time dominates each job, a per-call latency improvement
may have negligible job-level impact.

A useful deployment region therefore requires all of the following to some
degree: repeated state, scarce resources, exploitable advance information,
a supported actuator, and a benefit greater than control/interference costs.
The size of that region is an empirical question. It cannot be established by
showing that one regression predicts durations accurately.

\subsection{Comparison by responsibility, not headline speedup}
\begin{longtable}{@{}p{0.18\linewidth}p{0.35\linewidth}p{0.40\linewidth}@{}}
\caption{Relevant systems and the comparison ATFM needs. Descriptions refer to
cited mechanisms, not a claim that other systems cannot be extended.}\label{tab:landscape}\\
\toprule System & Established focus & What an ATFM comparison must isolate\\\midrule
\endfirsthead
\toprule System & Established focus & What an ATFM comparison must isolate\\\midrule
\endhead
vLLM & Engine scheduling, paged state, prefix reuse & Incremental benefit on a tuned engine with native reuse enabled\\
SGLang / HiCache & Program execution, radix reuse, cache tiers & Extra value over existing match/prefetch/eviction policies\\
TensorRT-LLM & Optimized execution, reuse and retention controls & Prediction benefit beyond available native priorities\\
LMCache & State retention, movement, cross-engine reuse & Which supported operation a forecast improves, including its cost\\
Sarathi-Serve & Chunked prefill and phase interference & Admission benefit after reasonable intra-engine scheduling\\
DistServe & Separate prefill/decode resource allocation & Whether forecast value survives transfer and pool queueing\\
Mooncake & KV-centric distributed serving and overload control & Extra session information versus existing cluster policy\\
Preble & Distributed prefix reuse and load balancing & Anticipatory timing versus locality/load decisions at arrival\\
Dynamo & Routing, transfer coordination, scaling & Session-conditional signals beyond aggregate predictors\\
llm-d & Kubernetes inference routing and deployment composition & Integration ownership and comparison with enabled policies\\
Continuum & Bounded retention and program scheduling & Conditional forecasts versus robust TTL decisions\\
ThunderAgent & Program-aware inference and tool resources & Additional forecast value at equal workflow visibility\\
Tutti / AsymCache & IO- or kernel-aware cache optimization & Complementarity with a more efficient cache data path\\
ATFM prototype & Tool signals, sampled forecasts, bounded directives & End-to-end benefit remains to be established on real workers\\\bottomrule
\end{longtable}

\subsection{What the repository currently demonstrates}
This assessment is anchored to the repository at \texttt{900d63b}, not an
assumption that documentation implies hardware validation. Local source paths
are listed in Appendix~\ref{sec:code-map}.

\begin{itemize}
\item Offline trace adapters, session forecasting, Monte Carlo aggregation,
calibration, and evaluation are implemented. Recorded H1/H1b findings depend
on corpus, split, horizon, and signal quality.
\item The simulator contains queues, modeled workers, cache eviction, admission,
and placement-related comparisons. It permits controlled counterfactuals but
is not a real inference engine.
\item The board, proxy, event ingestion, and directive paths have local HTTP and
mock-worker integration evidence. A recorded Mocker setup lacked required KV
capacity metrics and issued no holds or touches; that run verifies wiring,
not the efficacy of those controls.
\item The LMCache HTTP adapter is tested against controlled responses. Current
controller compatibility, completed transfers, and actual GPU cache effects
remain unverified on the target deployment.
\item Tier and replica-floor proposals can be logged; the default launcher is
not evidence of a fully connected production autoscaler.
\item The completed refactor passed 420 tests with five skips. Deterministic
policy and forecast fixtures retained output hashes. This supports code
maintenance, not a GPU goodput or model-quality claim.
\end{itemize}

The status document also states that a completed H100 performance study is not
part of the recorded local evidence. The strongest useful present claim is that
ATFM provides a testable forecasting/control prototype with measured CPU and
integration behavior. Production serving advantage is still a hypothesis.

\subsection{Known control costs and the Python/Rust decision}
The repository has already removed repeated full-history JSONL scans, batched
hold delivery, introduced persistent admission heaps, and reduced planner and
empirical-sampling overhead. These changes target work volume before changing
implementation language. The following bounds distinguish independent scale
variables; there is no single meaningful $N$ for the whole system.

\begin{table}[htbp]\centering\small
\caption{Selected ATFM costs after the documented control refactors. Bounds omit
model-specific setup and external IO unless named.}
\begin{tabularx}{\linewidth}{@{}lXX@{}}\toprule
Operation & Principal scaling & Remaining qualification\\\midrule
Incremental log drain & $O(\Delta E)$ newly read bytes/records & First read and restart still replay retained history\\
Heap admission updates & Amortized $O(\log Q)$ per heap operation & Due timers and compaction can cause bursts\\
Peer rank hints & $O(U)$ distinct-index scan & Expected $O(1)$ count updates; $U \leq Q$\\
Batched hold transport & $\lceil D/B\rceil$ RPCs, $O(D)$ payload & No whole-plan distributed transaction\\
Monte Carlo aggregation & Depends on $N,K,H$; typically at least $O(NK)$ sample work & Vectorization does not eliminate samples or output arrays\\
GDP slot search & Range pruning; worst-case linear in eligible slots per job & Resource minima can occur in different slots\\\bottomrule
\end{tabularx}\end{table}

Here $\Delta E$ is newly appended event data, $Q$ queued requests, $D$ directives,
$B$ directives per transport batch, $N$ sessions, $K$ samples, and $H$ horizons.
Worst-case GDP work can be $O(NF)$ for $F$ eligible slots; if both grow together,
this is quadratic. Additional sample and resource dimensions belong explicitly
in threshold construction. An index does not guarantee a logarithmic search
when its pruning condition is only necessary. By contrast, returning a full
$Q$-element diagnostic snapshot necessarily costs $\Omega(Q)$.

Recent local proxy profiling found repeated runtime-detection import work and
removed it by supplying a missing serving dependency. Its paired fake-worker
trials improved latency/CPU in that configuration, but prediction timeout rates
increased. That result is recorded in the proxy-profiling report. It demonstrates
why a lower CPU total is not sufficient evidence of a healthier control loop.
The September 28 implementation bounds unfinished proxy prediction jobs (four
by default), falls back immediately at capacity, and keeps abandoned running
jobs counted until completion. A monotonic deadline spans dispatch and execution;
board HTTP phases use its remaining budget instead of a 500\,ms minimum.
Caller outcomes and actual unfinished work are counted separately. This bounds
local work, not remote computation or hard real-time response. It can exchange
timeouts for immediate rejection, so timely prediction use must be measured
alongside request latency. At that revision, synchronous board computation and remaining linear
rank work remained investigation targets.


The subsequent board-isolation implementation moves ingestion, forecasting and
planning to one bounded worker and publishes complete scalar prediction views.
Built-in HTTP predictions use expected $O(1)$ lookup; a per-tick projection costs
$O(S+D)$ for $S$ sessions and $D$ duration observations visited across distinct
means. A monotonic age limit includes computation time and triggers fallback
when the last completed view is too old. Quantiles and directive JSON are
prepared on the worker; numerical algorithms and seeded draw order are unchanged.
Threads still share the GIL. In three fresh paired 1,024-session trials, timely
prediction use rises from 24.59\% to 31.61\% and board heartbeat p95 delay falls
from roughly 90\,ms to 1--2\,ms. However, client p95 increases in every pair,
including three further pairs with zero hold delay. This is evidence for board
responsiveness and availability, not a serving-latency improvement. The dated
board-isolation report preserves the regression, raw evidence and host limits.

Python is currently appropriate for orchestration and numerical experimentation,
with bulk numerical work handled by compiled libraries. Rust becomes attractive
for a \emph{measured} high-frequency CPU component when algorithmic fixes,
batching, and boundary costs have already been addressed. A native module should
have a narrow contract and differential tests; moving the whole system would
add maintenance cost and could leave the real bottleneck unchanged. Transferring
large KV tensors and scheduling small metadata are separate costs, so neither
``Python is always negligible'' nor ``Rust will fix saturation'' is justified.

\subsection{An evaluation that could validate or reject the project}
\paragraph{Stage 1: establish executable contracts.}
Pin model, tokenizer/template, engine, cache connector, transport, and GPU
configuration. Verify the cache identity, supported retention/movement
operation, completion acknowledgement, expiration behavior, and observed
residency. Record whether the action protects GPU, CPU, disk, or a remote copy.
A mock endpoint returning 200 is insufficient. Test disconnects, stale events,
partial failures, and fallback; measure whether a retry changes the same object
or accidentally creates duplicate work.

\paragraph{Stage 2: evaluate forecasts without leakage.}
Split by complete session or workflow family and time, not arbitrary events
from the same run. Report residual-time calibration and quantile loss at
operational horizons. Compare aggregate last-value/rate predictors, age-only
survival, tool-type models, and progress-conditioned models. Include missing
progress, censored calls, unseen tools, and correlated backend slowdowns.
A forecaster may be well calibrated overall and systematically wrong for the
largest, most expensive contexts; stratify accordingly.

For quantile $q_\alpha$ and outcome $y$, pinball loss is
\begin{equation}
 L_\alpha(y,q_\alpha)=(\alpha-\ind\{y<q_\alpha\})(y-q_\alpha).
\end{equation}
Report empirical interval coverage as well as width. Narrow incorrect intervals
are not better forecasts. Evaluate service demand jointly with timing when they
are dependent, and retain deterministic seeds for parity tests without treating
seeds as independent production workloads.

\paragraph{Stage 3: isolate decision value.}
Use the same tuned engine and cache configuration across arms: native policy;
robust TTL/working-set baseline; aggregate forecasting; session age-only
forecasting; progress-conditioned ATFM; and an oracle whose extra information is
explicit. Disable individual actuators to separate admission, retention,
prefetch, and scaling effects. A simulated ``ThunderAgent-style'' policy must
not be labelled a benchmark of the full ThunderAgent implementation. Likewise,
a clairvoyant cache schedule is an upper bound under its modeled constraints,
not a deployable competitor.

\paragraph{Stage 4: run closed-loop agent workloads on real workers.}
Vary sessions, prompt/output lengths, tool gaps, shared-prefix fraction, cache
capacity, transfer bandwidth, burst correlation, tenant mix, and forecast error.
Include low load where control can only add overhead, near-saturation load where
it may help, and sustained overload where admission must make tradeoffs. Use
both reproducible synthetic programs and real tool traces. Preserve causal
turn dependencies so latency changes alter future arrival times correctly.

\paragraph{Stage 5: account for the entire outcome.}
Measure offered/completed/rejected work, joint-SLO goodput, TTFT, ITL tails,
whole-job completion, fairness, cache bytes by tier, useful/wasted transfers,
recomputed tokens, GPU idle time, controller CPU, queueing, and prediction
deadline misses. Include warmup, shutdown/drain rules, retries, errors, and
load-generator lateness. Run paired workloads with alternating order and
confidence intervals at the independent workload/run level. Correlated tokens
from one request are not independent experimental replicates.

\paragraph{Success and falsification criteria.}
A positive result should show an improved user-relevant objective at equal
resource budget, no unreported quality loss, and acceptable fairness/overhead.
A negative result is informative if simple TTL or native policies match the
benefit, useful lead time is scarce, prediction deadlines are missed, or cache
movement dominates savings. The project remains useful only if the value of
extra information survives the full path from observation to completed action.


\paragraph{Measured proxy follow-up.}
The next implementation study separates prediction, admission and dispatch
waiting and replaces per-release peer-list construction with exact per-tier
index counts. A rank query scans $U_t$ distinct queued indices; updates are
expected $O(1)$ and storage is $O(U)$ across tiers. The worst case still has
$U_t=Q$, so this does not eliminate quadratic total rank work across a unique
backlog. Strict ties and the admission policy are preserved.
Six paired local trials reduced proxy CPU in every pair, while client p95
improved in four and regressed in two. At 1,024 sessions, the per-run p95 ranges
were 2.564--4.067\,s before and 1.823--2.630\,s after; variability precludes a
single reliable speedup claim. A separate larger-connection-pool experiment
regressed large-case latency and was reverted, despite opening fewer TCP
connections. All 46,080 unprofiled calls across the two studies succeeded;
one instrumented pool trial was incomplete and is identified separately.
The repository's dated proxy-ranking report preserves source revisions and raw
evidence. These measurements concern CPU and HTTP overhead with a fake worker,
not inference throughput, GPU memory movement, or production capacity.


\paragraph{Transport isolation and bounded pools.}
The subsequent local study isolates HTTP transport from the board and admission
queue, then balances unfinished response lifetimes across 16 smaller pools with
a combined 100-connection cap. Selection costs $O(P)$ for $P$ pools; internal
HTTP scans remain, and requests are not migrated between shards. The aggregate
idle cap rises from 20 to 100 with unchanged expiry. One shared verifying TLS
context and stock fallback for HTTP proxy discovery preserve operational choices.
The small-load fixture regressed latency, so automatic selection retains stock
HTTPX for initial admission windows of 20 or fewer; larger windows select shards.

All six full-proxy pairs improved client p95 and proxy CPU. At 1,024 sessions,
p95 ranges changed from 2.462--2.654\,s to 1.422--1.465\,s. Timely prediction use
fell from 35.50\% to 31.93\%, with substantially more caller timeouts, so improved
transport latency is not evidence of improved forecasting or control quality.
All 23,040 full-proxy and 82,944 formal component requests succeeded. These short
localhost nonstreaming trials exclude real inference and network/TLS costs;
real-socket lifecycle tests do not establish streaming performance. The dated
sharded-transport study retains measured revisions, unfavorable small-case
results, the final automatic-selection distinction and raw evidence.

\clearpage
